% claim.tex -- AI4Math claim of record (AI-generated, 2026-09-28; instantiated
% from harness/templates/claim/claim.tex per SOP 08b).
% Machine-readable theorem anchors are the "% LEAN:" lines; scripts/paper-lint.sh
% and the claim audit check them. All Lean listings are verbatim, byte-identical
% contiguous excerpts of frozen artifacts of tasks/20260927-mossad-cchord:
%   statements  : formalized/01-cchord-mods.lean (87 lines, sha256 a69d00e5...)
%   proof bodies: final/01-cchord-proved.lean    (262 lines, sha256 af39d35e...)
% Line ranges used: statement 22-42, 45-63, 66-86; final 43-67, 252-262.
% Verification: audit/listing-verify.txt (audit/inspect-tex.py --verify).
\documentclass[11pt]{article}

\usepackage[T1]{fontenc}
\usepackage[utf8]{inputenc}
\usepackage{amsmath,amssymb,amsthm}
\usepackage{graphicx}
\usepackage{hyperref}
\usepackage{xurl}
\setlength{\emergencystretch}{3em}
% lstlean.tex — Lean style for the listings package.
% Lineage: Jeremy Avigad (2015), adapted from Assia Mahboubi's SSR style;
% inspired by Olivier Verdier's unixode.sty for the Unicode literate table.
% No CTAN package or upstream repo exists; papers carry this file in their
% source package (cap set arXiv:1907.01449, sphere eversion arXiv:2210.07746).
% AI4Math adaptation (AI-generated, 2026-09-26): sorry/admit/sorryAx elevated
% to a dedicated error tier, matching the pipeline's 0-sorry constitution.
\usepackage{listings}
\usepackage{xcolor}

\definecolor{keywordcolor}{rgb}{0.7, 0.1, 0.1}   % red keywords
\definecolor{tacticcolor}{rgb}{0.0, 0.1, 0.6}    % blue tactics
\definecolor{commentcolor}{rgb}{0.4, 0.4, 0.4}   % grey comments
\definecolor{symbolcolor}{rgb}{0.0, 0.1, 0.6}    % blue symbols
\definecolor{sortcolor}{rgb}{0.1, 0.5, 0.1}      % green sorts
\definecolor{errorcolor}{rgb}{0.6, 0, 0}         % dark red: sorry/admit
\definecolor{stringcolor}{rgb}{0.5, 0.1, 0.5}    % purple strings

\lstdefinelanguage{lean}{
  % Anything between $ becomes LaTeX math mode
  mathescape=false,
  % Comments may or not include Latex commands
  texcl=false,
  % keywords, list taken from lean-syntax.txt
  morekeywords=[1]{import, prelude, protected, private, noncomputable,
    definition, meta, renaming, hiding, parameter, parameters, begin,
    constant, constants, lemma, variable, variables, theory, print,
    theorem, example, open, section, namespace, universe, universes, end,
    modifier, modifiers, using, export, exposing, axiom, inductive,
    coinductive, with, structure, class, instance, attribute, local, where,
    by, have, show, let, in, if, then, else, match, do, for, fun, assume,
    from, take, calc, include, omit, infix, infixl, infixr, notation,
    macro, syntax, elab, set_option, initialize, deriving, opaque, abbrev,
    mutual, partial, scoped, termination_by, decreasing_by},
  % Sorts
  morekeywords=[2]{Sort, Type, Prop},
  % tactics, list taken from lean-syntax.txt (tactic keywords)
  morekeywords=[3]{intro, intros, apply, exact, refine, rfl, simp, simp_all,
    simpa, simp_rw, rw, rewrite, erewrite, ring, ring_nf, omega, decide,
    norm_num, norm_cast, induction, cases, rcases, obtain, constructor,
    ext, funext, conv, congr, field_simp, linarith, nlinarith, positivity,
    tauto, trivial, aesop, use, by_cases, by_contra, push_neg, contrapose,
    symm, trans, assumption, contradiction, exfalso, specialize,
    generalize, revert, clear, rename_i, subst, unfold, delta, change,
    convert, split_ifs, interval_cases, fin_cases, zify, gcongr,
    nth_rewrite, suffices, generalizing, at, only, native_decide},
  % warnings - constitution tier: must never survive into a final proof
  morekeywords=[4]{sorry, admit, sorryAx},
  literate=
    {α}{{\ensuremath{\alpha}}}1
    {β}{{\ensuremath{\beta}}}1
    {γ}{{\ensuremath{\gamma}}}1
    {δ}{{\ensuremath{\delta}}}1
    {ε}{{\ensuremath{\varepsilon}}}1
    {ζ}{{\ensuremath{\zeta}}}1
    {η}{{\ensuremath{\eta}}}1
    {θ}{{\ensuremath{\theta}}}1
    {ι}{{\ensuremath{\iota}}}1
    {κ}{{\ensuremath{\kappa}}}1
    {λ}{{\ensuremath{\lambda}}}1
    {μ}{{\ensuremath{\mu}}}1
    {ν}{{\ensuremath{\nu}}}1
    {ξ}{{\ensuremath{\xi}}}1
    {π}{{\ensuremath{\pi}}}1
    {ρ}{{\ensuremath{\rho}}}1
    {σ}{{\ensuremath{\sigma}}}1
    {τ}{{\ensuremath{\tau}}}1
    {υ}{{\ensuremath{\upsilon}}}1
    {φ}{{\ensuremath{\varphi}}}1
    {χ}{{\ensuremath{\chi}}}1
    {ψ}{{\ensuremath{\psi}}}1
    {ω}{{\ensuremath{\omega}}}1
    {Γ}{{\ensuremath{\Gamma}}}1
    {Δ}{{\ensuremath{\Delta}}}1
    {Θ}{{\ensuremath{\Theta}}}1
    {Λ}{{\ensuremath{\Lambda}}}1
    {Ξ}{{\ensuremath{\Xi}}}1
    {Π}{{\ensuremath{\Pi}}}1
    {Σ}{{\ensuremath{\Sigma}}}1
    {Φ}{{\ensuremath{\Phi}}}1
    {Ψ}{{\ensuremath{\Psi}}}1
    {Ω}{{\ensuremath{\Omega}}}1
    {ℵ}{{\ensuremath{\aleph}}}1
    {∀}{{\ensuremath{\forall}}}1
    {∃!}{{\ensuremath{\exists}!}}1
    {∃}{{\ensuremath{\exists}}}1
    {∈}{{\ensuremath{\in}}}1
    {∉}{{\ensuremath{\notin}}}1
    {∘}{{\ensuremath{\circ}}}1
    {∩}{{\ensuremath{\cap}}}1
    {∪}{{\ensuremath{\cup}}}1
    {⊆}{{\ensuremath{\subseteq}}}1
    {⊂}{{\ensuremath{\subset}}}1
    {⊇}{{\ensuremath{\supseteq}}}1
    {⊃}{{\ensuremath{\supset}}}1
    {⊄}{{\ensuremath{\nsubseteq}}}1
    {∅}{{\ensuremath{\emptyset}}}1
    {∧}{{\ensuremath{\wedge}}}1
    {∨}{{\ensuremath{\vee}}}1
    {¬}{{\ensuremath{\neg}}}1
    {⊤}{{\ensuremath{\top}}}1
    {⊥}{{\ensuremath{\bot}}}1
    {↔}{{\ensuremath{\leftrightarrow}}}1
    {→}{{\ensuremath{\rightarrow}}}1
    {←}{{\ensuremath{\leftarrow}}}1
    {↑}{{\ensuremath{\uparrow}}}1
    {⇒}{{\ensuremath{\Rightarrow}}}1
    {⇐}{{\ensuremath{\Leftarrow}}}1
    {⟹}{{\ensuremath{\Longrightarrow}}}1
    {↦}{{\ensuremath{\mapsto}}}1
    {≤}{{\ensuremath{\leq}}}1
    {≥}{{\ensuremath{\geq}}}1
    {≠}{{\ensuremath{\neq}}}1
    {≈}{{\ensuremath{\approx}}}1
    {≡}{{\ensuremath{\equiv}}}1
    {≃}{{\ensuremath{\simeq}}}1
    {≅}{{\ensuremath{\cong}}}1
    {×}{{\ensuremath{\times}}}1
    {·}{{\ensuremath{\cdot}}}1
    {±}{{\ensuremath{\pm}}}1
    {⊕}{{\ensuremath{\oplus}}}1
    {⊗}{{\ensuremath{\otimes}}}1
    {⋆}{{\ensuremath{\star}}}1
    {▸}{{\ensuremath{\blacktriangleright}}}1
    {⟨}{{\ensuremath{\langle}}}1
    {⟩}{{\ensuremath{\rangle}}}1
    {⦃}{{\ensuremath{\{\!\mid}}}1
    {⦄}{{\ensuremath{\mid\!\}}}}1
    {⟦}{{\ensuremath{[\![}}}1
    {⟧}{{\ensuremath{]\!]}}}1
    {⌊}{{\ensuremath{\lfloor}}}1
    {⌋}{{\ensuremath{\rfloor}}}1
    {⌈}{{\ensuremath{\lceil}}}1
    {⌉}{{\ensuremath{\rceil}}}1
    {√}{{\ensuremath{\surd}}}1
    {∣}{{\ensuremath{\mid}}}1
    {∤}{{\ensuremath{\nmid}}}1
    {∎}{{\ensuremath{\blacksquare}}}1
    {⋯}{{\ensuremath{\cdots}}}1
    {…}{{\ensuremath{\ldots}}}1
    {∑}{{\ensuremath{\sum}}}1
    {∏}{{\ensuremath{\prod}}}1
    {⁻¹}{{\ensuremath{^{-1}}}}1
    {¹}{{\ensuremath{^1}}}1
    {²}{{\ensuremath{^2}}}1
    {³}{{\ensuremath{^3}}}1
    {ⁿ}{{\ensuremath{^n}}}1
    {ᵐ}{{\ensuremath{^m}}}1
    {₀}{{\ensuremath{_0}}}1
    {₁}{{\ensuremath{_1}}}1
    {₂}{{\ensuremath{_2}}}1
    {₃}{{\ensuremath{_3}}}1
    {₄}{{\ensuremath{_4}}}1
    {₅}{{\ensuremath{_5}}}1
    {ₙ}{{\ensuremath{_n}}}1
    {ₘ}{{\ensuremath{_m}}}1
    {ₖ}{{\ensuremath{_k}}}1
    {ᵢ}{{\ensuremath{_i}}}1
    {ⱼ}{{\ensuremath{_j}}}1
    {ℤ}{{\ensuremath{\mathbb{Z}}}}1
    {ℕ}{{\ensuremath{\mathbb{N}}}}1
    {ℚ}{{\ensuremath{\mathbb{Q}}}}1
    {ℝ}{{\ensuremath{\mathbb{R}}}}1
    {ℂ}{{\ensuremath{\mathbb{C}}}}1
    {𝔽}{{\ensuremath{\mathbb{F}}}}1
    {𝒫}{{\ensuremath{\mathcal{P}}}}1
    {∗}{{\ensuremath{\ast}}}1
    {•}{{\ensuremath{\bullet}}}1,
  % Comments
  morecomment=[l]{--},
  morecomment=[s]{/-}{-/},
  morestring=[b]",
  stringstyle=\color{stringcolor},
  keepspaces=true,
  keywordstyle=[1]\color{keywordcolor}\bfseries,
  keywordstyle=[2]\color{sortcolor}\bfseries,
  keywordstyle=[3]\color{tacticcolor}\bfseries,
  keywordstyle=[4]\color{errorcolor}\bfseries\underline,
  escapeinside={(*@}{@*)},
  columns=fullflexible,
  basicstyle=\ttfamily\small,
  commentstyle=\color{commentcolor}\itshape,
  breaklines=true,
  showstringspaces=false,
  frame=single,
  framesep=2mm,
  xleftmargin=4mm,
  numbers=left,
  numberstyle=\tiny\color{commentcolor},
  numbersep=6pt,
}

% Inline Lean code in prose: \lean{Int.ModEq.pow}
\newcommand{\lean}[1]{\lstinline[language=lean,basicstyle=\ttfamily\normalsize]|#1|}

\lstset{language=lean}

% --- local rendering patch (2026-09-28) ---
% tectonic/XeTeX does not apply the listings literate table to multi-byte
% characters, so the Unicode set used by the listings of THIS note is mapped to
% math glyphs as active characters via the standard \lccode/\lowercase idiom
% (ASCII-only source, no extra packages). Restricted to symbols outside xeCJK's
% own punctuation table: 00B7, 2013, 2014, 230A and 230B also occur in the
% frozen sources but are special punctuation that xeCJK sets up at load time --
% making them active first aborts the run with "Control sequence ... already
% defined" (observed 2026-09-27); those five route through xeCJK + SimSun
% instead. Rendering only: listing source bytes are untouched, so the
% byte-exactness checks are unaffected, and the \ifdefined guard leaves the
% pdflatex path identical to the template.
\ifdefined\XeTeXversion
  {\lccode`\~="2022 \lowercase{\global\catcode`~=\active \global\def~{\ensuremath{\bullet}}}}
  {\lccode`\~="2115 \lowercase{\global\catcode`~=\active \global\def~{\ensuremath{\mathbb{N}}}}}
  {\lccode`\~="2124 \lowercase{\global\catcode`~=\active \global\def~{\ensuremath{\mathbb{Z}}}}}
  {\lccode`\~="2190 \lowercase{\global\catcode`~=\active \global\def~{\ensuremath{\leftarrow}}}}
  {\lccode`\~="2192 \lowercase{\global\catcode`~=\active \global\def~{\ensuremath{\rightarrow}}}}
  {\lccode`\~="2194 \lowercase{\global\catcode`~=\active \global\def~{\ensuremath{\leftrightarrow}}}}
  {\lccode`\~="2200 \lowercase{\global\catcode`~=\active \global\def~{\ensuremath{\forall}}}}
  {\lccode`\~="2203 \lowercase{\global\catcode`~=\active \global\def~{\ensuremath{\exists}}}}
  {\lccode`\~="2227 \lowercase{\global\catcode`~=\active \global\def~{\ensuremath{\wedge}}}}
  {\lccode`\~="2228 \lowercase{\global\catcode`~=\active \global\def~{\ensuremath{\vee}}}}
  {\lccode`\~="2261 \lowercase{\global\catcode`~=\active \global\def~{\ensuremath{\equiv}}}}
  {\lccode`\~="2264 \lowercase{\global\catcode`~=\active \global\def~{\ensuremath{\leq}}}}
  {\lccode`\~="2265 \lowercase{\global\catcode`~=\active \global\def~{\ensuremath{\geq}}}}
  {\lccode`\~="22A2 \lowercase{\global\catcode`~=\active \global\def~{\ensuremath{\vdash}}}}
  {\lccode`\~="27E8 \lowercase{\global\catcode`~=\active \global\def~{\ensuremath{\langle}}}}
  {\lccode`\~="27E9 \lowercase{\global\catcode`~=\active \global\def~{\ensuremath{\rangle}}}}
  {\lccode`\~="00AC \lowercase{\global\catcode`~=\active \global\def~{\ensuremath{\neg}}}}
  {\lccode`\~="2212 \lowercase{\global\catcode`~=\active \global\def~{\ensuremath{-}}}}
  {\lccode`\~="2223 \lowercase{\global\catcode`~=\active \global\def~{\ensuremath{\mid}}}}
\fi
\ifdefined\XeTeXversion
  \usepackage{xeCJK}
  \setCJKmainfont{SimSun}
\fi

\newtheorem{theorem}{Theorem}
\theoremstyle{definition}
\newtheorem{conjecture}{Conjecture}

\title{Mod-4 congruence and parity laws for the residue-class recurrences of
chorded-cycle matching counts\\ (Lean 4 machine-checked claim of record)}
\author{Aurora0134\thanks{Contact: via the GitHub account Aurora0134; ORCID to
be added at Zenodo upload.}}
\date{September 28, 2026}

\begin{document}
\maketitle

\begin{abstract}
Let $a(n)$ be the number of matchings (including the empty one) of the graph on
$\mathbb{Z}/n\mathbb{Z}$ with the cycle edges plus all chords
$\{i,i+\lfloor n/3\rfloor\}$. Exact enumeration for $n\ge 3$ starts
$4,7,11,51,99,191,382,780,\dots$; the sequence splits by $n\bmod 3$ into three
subfamilies whose terms are reproduced, up to $n=150$, by integer recurrences of
orders $6$, $8$ and $8$ fitted over $\mathbb{Q}$ with Berlekamp--Massey. We state
and machine-check, in Lean 4 (v4.34.0, mathlib v4.34.0), three congruence laws
for those recurrence-defined sequences: constant $3\bmod 4$; a two-valued
$\bmod 4$ law along $j+1\bmod 3$ for $j\ge 1$; and parity characterised by
$(j+1)\bmod 5=3$. Every proof compiles with no \texttt{sorry} and
\texttt{\#print axioms} reports only \{propext, Quot.sound, Classical.choice\}.
Novelty adjudication: no occupying record found in OEIS, zbMATH, Math Stack
Exchange, arXiv or the public Lean ecosystem under the query discipline recorded
in the deposit; value tier \emph{new sequence / new recurrence}, not new
mathematics. The combinatorial identification is not kernel-proved and is stated
as a conjecture. This note is a priority claim of record, not a full paper.
\end{abstract}

\section{Introduction}

Three theorems about three integer sequences defined by initial values and
linear recurrences are reported here; the sequences come from counting matchings
of cycles augmented by chords of length $\lfloor n/3\rfloor$, and the recurrences
were obtained by exact rational fitting of enumerated data
\cite{massey,godsil-gutman}. The proofs were produced by the AI4Math pipeline, in
which LLM workers generate statements and proof scripts and the Lean 4 kernel is
the sole adjudicator at every gate: statements are frozen by hash before
proving, and a separate audit re-computes the symmetric difference, scans the
whole file for \texttt{sorry}/\texttt{admit} including comments, and re-prints
the axiom dependencies \cite{lean4,mathlib}. Verification strength is stated
below verbatim. This note is a priority claim of record deposited so that a
dated, citable public record exists ahead of any narrative paper; the fuller
paper for the same three theorems is prepared separately and will reference
this deposit. No literature review is attempted here.

\section{Statement}\label{sec:statement}

Throughout, the three sequences are \emph{defined} by the displayed initial
values and recurrences; $0$-based indexing, with $c_r(j)$ corresponding to the
combinatorial $a(3(j+1)+r)$.

% LEAN: tasks/20260927-mossad-cchord :: cchordR2_mod_four grade=完全证明
\begin{theorem}[constant congruence]\label{thm:r2}
For the sequence $c_2$ of Listing~\ref{lst:r2} (nine seeds, order-$8$
recurrence), $c_2(k)\equiv 3\pmod 4$ for every $k\in\mathbb{N}$.
\end{theorem}

\begin{lstlisting}[caption={Frozen formal statement, verbatim lines 22--42 of
\texttt{formalized/01-cchord-mods.lean} (sha256 \texttt{a69d00e5}).},
label={lst:r2}]
/-- **序列 cchordR2（r=2 子族，阶 8 递推）**。
初值 c(0..8) = 11, 191, 1475, 10615, 77623, 565895, 4127167, 30100411, 219522983
（= counts.txt 中 a(5),a(8),a(11),a(14),a(17),a(20),a(23),a(26),a(29)）；
递推 c(n+9) = 4c(n+8)+22c(n+7)+18c(n+6)−22c(n+5)−16c(n+4)+10c(n+3)+2c(n+2)−c(n+1)。 -/
def cchordR2 : ℕ → ℤ
  | 0 => 11
  | 1 => 191
  | 2 => 1475
  | 3 => 10615
  | 4 => 77623
  | 5 => 565895
  | 6 => 4127167
  | 7 => 30100411
  | 8 => 219522983
  | n + 9 =>
      4 * cchordR2 (n + 8) + 22 * cchordR2 (n + 7) + 18 * cchordR2 (n + 6)
        - 22 * cchordR2 (n + 5) - 16 * cchordR2 (n + 4) + 10 * cchordR2 (n + 3)
        + 2 * cchordR2 (n + 2) - cchordR2 (n + 1)

/-- **C1（主攻）**：r=2 子族 mod 4 恒等于 3。 -/
theorem cchordR2_mod_four : ∀ k : ℕ, cchordR2 k ≡ 3 [ZMOD 4] := by
\end{lstlisting}

% LEAN: tasks/20260927-mossad-cchord :: cchordR0_mod_four grade=完全证明
\begin{theorem}[two-valued congruence along $j+1\bmod 3$]\label{thm:r0}
For the sequence $c_0$ of Listing~\ref{lst:r0} (seven seeds, order-$6$
recurrence) and every $j\ge 1$,
$c_0(j)\equiv 2\pmod 4$ if $3\mid(j+1)$ and $c_0(j)\equiv 3\pmod 4$ otherwise.
The hypothesis is needed: $c_0(0)=4\equiv 0\pmod 4$.
\end{theorem}

\begin{lstlisting}[caption={Frozen formal statement, verbatim lines 45--63 of
\texttt{formalized/01-cchord-mods.lean}.},label={lst:r0}]
/-- **序列 cchordR0（r=0 子族，阶 6 递推）**。
初值 c(0..6) = 4, 51, 382, 2743, 19907, 144054, 1043147
（= counts.txt 中 a(3),a(6),a(9),a(12),a(15),a(18),a(21)）；
递推 c(n+7) = 5c(n+6)+16c(n+5)+3c(n+4)−10c(n+3)−c(n+2)+c(n+1)。 -/
def cchordR0 : ℕ → ℤ
  | 0 => 4
  | 1 => 51
  | 2 => 382
  | 3 => 2743
  | 4 => 19907
  | 5 => 144054
  | 6 => 1043147
  | n + 7 =>
      5 * cchordR0 (n + 6) + 16 * cchordR0 (n + 5) + 3 * cchordR0 (n + 4)
        - 10 * cchordR0 (n + 3) - cchordR0 (n + 2) + cchordR0 (n + 1)

/-- **C2（伴随）**：r=0 子族 mod 4 三分支（j≥1 起；j=0 处 ≡0 为已知首例外）。 -/
theorem cchordR0_mod_four (j : ℕ) (hj : 1 ≤ j) :
    cchordR0 j ≡ (if 3 ∣ (j + 1) then 2 else 3) [ZMOD 4] := by
\end{lstlisting}

% LEAN: tasks/20260927-mossad-cchord :: cchordR1_even_iff grade=完全证明
\begin{theorem}[parity characterisation]\label{thm:r1}
For the sequence $c_1$ of Listing~\ref{lst:r1} (nine seeds, order-$8$
recurrence) and every $j\in\mathbb{N}$,
$2\mid c_1(j)$ if and only if $(j+1)\bmod 5=3$.
\end{theorem}

\begin{lstlisting}[caption={Frozen formal statement, verbatim lines 66--86 of
\texttt{formalized/01-cchord-mods.lean}.},label={lst:r1}]
/-- **序列 cchordR1（r=1 子族，阶 8 递推）**。
初值 c(0..8) = 7, 99, 780, 5591, 41281, 302139, 2217121, 16257026, 119231297
（= counts.txt 中 a(4),a(7),a(10),a(13),a(16),a(19),a(22),a(25),a(28)）；
递推 c(n+9) = 4c(n+8)+23c(n+7)+14c(n+6)−23c(n+5)−14c(n+4)+9c(n+3)+2c(n+2)−c(n+1)。 -/
def cchordR1 : ℕ → ℤ
  | 0 => 7
  | 1 => 99
  | 2 => 780
  | 3 => 5591
  | 4 => 41281
  | 5 => 302139
  | 6 => 2217121
  | 7 => 16257026
  | 8 => 119231297
  | n + 9 =>
      4 * cchordR1 (n + 8) + 23 * cchordR1 (n + 7) + 14 * cchordR1 (n + 6)
        - 23 * cchordR1 (n + 5) - 14 * cchordR1 (n + 4) + 9 * cchordR1 (n + 3)
        + 2 * cchordR1 (n + 2) - cchordR1 (n + 1)

/-- **C3（伴随）**：r=1 子族奇偶 iff——偶数当且仅当 (j+1) % 5 = 3。 -/
theorem cchordR1_even_iff (j : ℕ) : 2 ∣ cchordR1 j ↔ (j + 1) % 5 = 3 := by
\end{lstlisting}

\begin{conjecture}[combinatorial bridge, not proved here]\label{conj:bridge}
For $r\in\{0,1,2\}$ and all $j\ge 0$, $c_r(j)=a(3(j+1)+r)$, where $a$ counts
matchings of the chorded cycle described in the abstract.
\end{conjecture}

\section{Proof}\label{sec:proof}

All three proofs are strong induction on the index with the seed window
discharged by computation and the step closed by congruence algebra inside
\texttt{Int.ModEq}; Theorems~\ref{thm:r0} and \ref{thm:r1} add a case split on
$(j+1)\bmod 3$ and $(j+1)\bmod 5$ respectively. The complete proof of
Theorem~\ref{thm:r2} is reproduced verbatim below. The other two declarations
are longer --- Theorem~\ref{thm:r0} spans lines 86--137 and Theorem~\ref{thm:r1}
lines 159--262 of the same final file --- and are excerpted here only where the
argument has a distinctive shape: the closing assembly of Theorem~\ref{thm:r1}
from its modulo-$2$ strengthening, which is where the if-and-only-if is won. The
complete $262$-line file ships in this deposit under \texttt{proofs/}.

\begin{lstlisting}[caption={Complete proof of Theorem~\ref{thm:r2}, verbatim
lines 43--67 of \texttt{final/01-cchord-proved.lean} (sha256
\texttt{af39d35e}).},label={lst:p2}]
theorem cchordR2_mod_four : ∀ k : ℕ, cchordR2 k ≡ 3 [ZMOD 4] := by
  intro n
  induction n using Nat.strong_induction_on with
  | _ n ih =>
    rcases lt_or_ge n 9 with h | h
    · interval_cases n <;> decide
    · obtain ⟨j, rfl⟩ : ∃ j, n = j + 9 := ⟨n - 9, by omega⟩
      have rec9 : ∀ m : ℕ, cchordR2 (m + 9)
          = 4 * cchordR2 (m + 8) + 22 * cchordR2 (m + 7) + 18 * cchordR2 (m + 6)
            - 22 * cchordR2 (m + 5) - 16 * cchordR2 (m + 4) + 10 * cchordR2 (m + 3)
            + 2 * cchordR2 (m + 2) - cchordR2 (m + 1) := fun m => by
        simp only [cchordR2]
      rw [rec9 j]
      have h8 : cchordR2 (j + 8) ≡ 3 [ZMOD 4] := ih (j + 8) (by omega)
      have h7 : cchordR2 (j + 7) ≡ 3 [ZMOD 4] := ih (j + 7) (by omega)
      have h6 : cchordR2 (j + 6) ≡ 3 [ZMOD 4] := ih (j + 6) (by omega)
      have h5 : cchordR2 (j + 5) ≡ 3 [ZMOD 4] := ih (j + 5) (by omega)
      have h4 : cchordR2 (j + 4) ≡ 3 [ZMOD 4] := ih (j + 4) (by omega)
      have h3 : cchordR2 (j + 3) ≡ 3 [ZMOD 4] := ih (j + 3) (by omega)
      have h2 : cchordR2 (j + 2) ≡ 3 [ZMOD 4] := ih (j + 2) (by omega)
      have h1 : cchordR2 (j + 1) ≡ 3 [ZMOD 4] := ih (j + 1) (by omega)
      have key := (((((((h8.mul_left 4).add (h7.mul_left 22)).add (h6.mul_left 18)).sub
          (h5.mul_left 22)).sub (h4.mul_left 16)).add (h3.mul_left 10)).add
          (h2.mul_left 2)).sub h1
      exact key.trans (by decide)
\end{lstlisting}

\begin{lstlisting}[caption={Excerpt: assembly of Theorem~\ref{thm:r1} from the
modulo-$2$ strengthening \texttt{main} proved at lines 160--251 of the same
file; verbatim lines 252--262 (contiguous block; the complete proofs of
Theorem~\ref{thm:r0} and Theorem~\ref{thm:r1} are lines 86--137 and 159--262,
in \texttt{proofs/}).},label={lst:p3}]
  have h := main j
  rw [← Int.modEq_zero_iff_dvd]
  constructor
  · intro hd
    by_cases hc : (j + 1) % 5 = 3
    · exact hc
    · rw [if_neg hc] at h
      exact absurd (h.symm.trans hd) (by decide)
  · intro hc
    rw [if_pos hc] at h
    exact h
\end{lstlisting}

\section{Verification evidence}

\begin{itemize}
\item Toolchain: Lean \texttt{v4.34.0} (\texttt{leanprover/lean4:v4.34.0}),
mathlib4 \texttt{v4.34.0}, revision \texttt{5ed29652}. Reproduction: with that
pin, \texttt{lake env lean} on \texttt{proofs/01-cchord-proved.lean} compiles
the file with exit code $0$ (measured $29.1$\,s for the strict pass).
\item Statement integrity: the frozen statement file
\texttt{formalized/01-cchord-mods.lean} ($87$ lines) has sha256
\texttt{a69d00e537675c40bd13e82344004e7f}\allowbreak
\texttt{1abd370dc3d9e543396fbf1818c3a03b}; the
final file \texttt{final/01-cchord-proved.lean} ($262$ lines) has sha256
\texttt{af39d35ea29c9a01369473140319f9c3}\allowbreak
\texttt{30d326601463839875d11b14bbba2237}.
The final file was compared with the frozen file by symmetric difference over
the statement region: zero drift; the only differences are the three proof
bodies replacing the \texttt{sorry} placeholders and one header-comment line.
\item Kernel check: compiles with $0$ \texttt{sorry} (whole-file scan including
comments, case-insensitive, also for \texttt{admit}/\texttt{sorryAx});
\texttt{\#print axioms} output, verbatim (kept literal; the smaller monospace
width is typesetting only):
\begin{lstlisting}[basicstyle=\ttfamily\footnotesize,frame=none]
'cchordR2_mod_four' depends on axioms: [propext, Classical.choice, Quot.sound]
'cchordR0_mod_four' depends on axioms: [propext, Classical.choice, Quot.sound]
'cchordR1_even_iff' depends on axioms: [propext, Classical.choice, Quot.sound]
\end{lstlisting}
The strict re-compile emits $68$ deprecation warnings
(\texttt{if\_neg}/\texttt{if\_pos}) and no error; warnings are not treated as a
refusal ground by the audit, and are recorded here. The log ships as
\texttt{audit/audit-strict-compile.log}, with machine-specific path prefixes
replaced by \texttt{<repo>/} for the public deposit.
\item Grading by the audit department: all three theorems
\emph{complete proof}; lowest grade reached for the case = \emph{complete
proof}; no scaffolding route was used.
\item Audit chain (originals in this deposit under \texttt{audit/}): final
adjudication \texttt{final-01-cchord.txt} (four independent checks),
well-definedness verdict \texttt{welldef-verdict.md} with gate logs
\texttt{gate-c1.txt}, \texttt{gate-c2.txt}, \texttt{gate-c3.txt} and the
rechecker log \texttt{gate-recheck-frozen.txt}, axiom printout
\texttt{audit-axioms-sidecar.log} (the three lines quoted above are its whole
content), and this note's claim check \texttt{claim-check.md} --- added to
\texttt{audit/} after the audit department returns its adjudication, and a
prerequisite for publication. The semantic-fidelity check (roundtrip
re-translation) of the underlying task ran on the same model node as its
formalisation; that independence weakness is mitigated, not removed, by kernel
adjudication and manual review.
\end{itemize}

\section{Novelty evidence}\label{sec:novelty}

Adjudication copied from the pipeline's exclusivity review (C9 gate), which ran
before the problem was accepted and re-derived the data independently: verdict
\textbf{no occupying record found}, value tier \textbf{new sequence / new
recurrence}. It is not a claim of new mathematics and not a first-discovery
claim. Summary by channel; the full query log, including every zero-hit query,
is in this deposit at \texttt{audit/c9-record.md}.

\begin{center}
\resizebox{\textwidth}{!}{%
\begin{tabular}{lll}
\hline
channel & queries & result \\
\hline
OEIS sequence search \cite{oeis} & 9 & zero hits: supplied seed and its
drop-first / drop-two variants (3), real sequence and its drop-first /
drop-two variants (3), $\times2$ (1), $+1$ (1), tail segment (1) \\
OEIS per-residue subsequences & 3 & zero hits (\texttt{4,51,382,2743},
\texttt{7,99,780,5591}, \texttt{11,191,1475,10615}) \\
OEIS perfect-matching variant & 2 & zero hits \\
OEIS keywords & 2 & 10 hits each, none matching counts of this family \\
zbMATH & 4 & \texttt{Hosoya index circulant graph}: clean zero reading; three
other queries return hits in other directions (Pfaffian property; pancyclicity
and Hamilton-type chorded-cycle results \cite{chorded-survey}; matching
polynomials) \\
Math Stack Exchange, arXiv & 2 & zero hits \\
OpenAlex & 1 & HTTP 429 on this candidate, not verified (gap disclosed) \\
public Lean ecosystem & 2 & zero hits; one repository tree walk rate-limited \\
library layer & 3 & mathlib has \texttt{circulantGraph} \cite{mathlib-circulant}
but no matching-count theorem; compfiles none; sequencelib none \\
\hline
\end{tabular}%
}
\end{center}

Two points of the record are part of the evidence. The externally supplied seed
began $2,2,\dots$, and those two terms are artefacts of that pipeline's
degenerate $n=1,2$ cases; the sequence used here starts at $n=3$ with $4$. And
the general web search layer was unavailable during the review, so literature
coverage is zbMATH-, arXiv- and Math-Stack-Exchange-centric; a
Chinese-language database was not searched. If any of these channels later
returns a hit on chorded-cycle matching counts, the tier above weakens
accordingly, and this deposit will be amended by a later version rather than by
editing this record.

\section{Scope and limitations}

Grade and boundary, stated as the audit states them. (a) The theorems cover the
recurrence-defined sequences only; their equality with matching counts of the
chorded cycle is not kernel-proved and is Conjecture~\ref{conj:bridge}. Its
empirical support is enumeration by exact programs: the counting main program
produced $3\le n\le 150$, cross-checked against a brute-force enumerator for
$3\le n\le 16$ and against a second layout for $3\le n\le 36$, and against the
externally supplied seed only for $3\le n\le 15$; the enumerators themselves are
not verified, and the scripts and their outputs stay in the working repository
under \texttt{tasks/20260927-mossad-cchord/}. (b)
Theorem~\ref{thm:r0} holds for $j\ge 1$; at $j=0$ the value is $4$, which is
$\equiv 0\pmod 4$. (c) Indexing is $0$-based: $c_r(j)$ corresponds to
$a(3(j+1)+r)$. (d) The recurrences are computational findings: Berlekamp--Massey
fits over $\mathbb{Q}$ verified digit-for-digit against enumeration, not
derived from a transfer-matrix argument, and the natural uniform step-$3$
ansatz is refuted on the enumerated window by exact rational arithmetic ($136$
failures out of $139$ test indices, $12\le n\le 150$) --- a computation, not a
theorem of this note. (e) The gate-four sign-off of the underlying task was
exercised by the author's instruction to deposit this claim on 2026-09-27, while
the task report file still carries its draft header; both facts are recorded
here as they stand. (f) This note carries no literature review and is a deposit
record, not a survey; the narrative paper for the same result is separate.

\section*{Data availability}

The deposit package for this claim contains the claim note (LaTeX source and
PDF), the Lean sources (\texttt{proofs/}: frozen statement file, final proof
file, \texttt{lean-toolchain}), the verification and audit logs, and the full
C9 query log with zero-hit entries (\texttt{audit/}). Zenodo: the DOI is
reserved at upload by the author and is recorded in \texttt{metadata/README.md}
of this package once created; see \texttt{UPLOAD.md} in the source repository.
Public mirror of this deposit: the claim lane's repository
\url{https://github.com/Aurora0134/ai4math-claims}, directory
\texttt{chorded-cycle-mod4/} (pushed under the author's standing instruction;
the commit and timestamp are recorded in the source repository's
\texttt{claims/chorded-cycle-mod4/card.md}).
Claim note under CC BY 4.0, code under MIT.

\section*{Use of generative AI}

(a) AI performed: candidate generation and the exclusivity review,
autoformalisation of the three statements, the roundtrip semantic check, proof
search and proof-script writing, error classification and repair, and the
assembly of this note including its verbatim listings. (b) Model, node and
budget: all task-stage sampling ran on one node, wire-id
\texttt{anthropic/a6api-main/kimi-k3}, with no cross-node switching, and
consumed $1$ LLM call of a permitted $8$ and $15$ kernel compiles of a
permitted $24$, copied verbatim from the task ledger. This note consumed $0$
pipeline LLM calls against a cap of $4$; the drafting session ran with its
model selection key in auto-routing mode (last observed relay node
\texttt{stepfun/step-5-preview}), which is disclosed rather than replaced by the
task node's name. (c) Final adjudication: all statements and proofs reported
here were checked by the Lean 4 kernel: every proof compiles with no
\texttt{sorry}, and \texttt{\#print axioms} reports only \{propext, Quot.sound,
Classical.choice\}. (d) The AI system is not an author; the human author
reviewed the evidence and is responsible for all content.

\section*{Author contributions}

CRediT-style: the human author adjudicated topic selection, statement semantics
and the deposit decision; the AI4Math pipeline (an LLM) generated the candidate
propositions, the proof searches and the text drafts. The AI system is not
listed as an author.

\begin{thebibliography}{9}
% Reused from papers/chorded-cycle-mod4/paper.tex, where each entry was verified
% on 2026-09-27 against Crossref title metadata (four DOIs), the arXiv export API
% plus a 200 landing page (preprint), and 200 fetches for oeis.org and the pinned
% mathlib file. Raw responses: papers/chorded-cycle-mod4/audit/bib-verify/. No
% entry is quoted from memory; the DOI carried by our paper template for
% reference lean4 was found to be wrong (it resolves to a Collatz paper in the
% same volume) and is corrected to ..._37 here, verified by title metadata.
\bibitem{lean4}
L.~de Moura and S.~Ullrich.
\newblock The Lean 4 theorem prover and programming language.
\newblock \emph{Automated Deduction -- CADE-28}, LNCS, pp.~625--635, 2021.
\url{https://doi.org/10.1007/978-3-030-79876-5_37}

\bibitem{mathlib}
The mathlib Community.
\newblock The Lean mathematical library.
\newblock \emph{CPP 2020}, pp.~367--381, 2020.
\url{https://doi.org/10.1145/3372885.3373824}

\bibitem{mathlib-circulant}
I.~Renison and B.~Mehta.
\newblock \texttt{Mathlib/Combinatorics/SimpleGraph/Circulant.lean}, mathlib4,
tag \texttt{v4.34.0}.
\newblock
\url{https://raw.githubusercontent.com/leanprover-community/mathlib4/v4.34.0/Mathlib/Combinatorics/SimpleGraph/Circulant.lean}

\bibitem{godsil-gutman}
C.~D. Godsil and I.~Gutman.
\newblock On the theory of the matching polynomial.
\newblock \emph{Journal of Graph Theory} 5 (1981), pp.~137--144.
\url{https://doi.org/10.1002/jgt.3190050203}

\bibitem{massey}
J.~L. Massey.
\newblock Shift-register synthesis and BCH decoding.
\newblock \emph{IEEE Transactions on Information Theory} 15 (1969),
pp.~122--127.
\url{https://doi.org/10.1109/TIT.1969.1054260}

\bibitem{chorded-survey}
B.~Elliott, R.~Gould and K.~Hirohata.
\newblock On degree sum conditions and vertex-disjoint chorded cycles.
\newblock \emph{arXiv:1911.08686 [math.CO]}, 2019.
\url{https://arxiv.org/abs/1911.08686}

\bibitem{oeis}
The OEIS Foundation.
\newblock \emph{The On-Line Encyclopedia of Integer Sequences}.
\newblock \url{https://oeis.org/}
\end{thebibliography}

\end{document}
